\documentclass[main]{subfiles}


\title[Evaluation and Benchmarking]{Evaluation and Benchmarking}
\subtitle{The Big Picture}
\author[Marcel Wever]{Marcel Wever}
\date{}

\titlebackground*{images/AutoML_ML_Evaluation}

%\institute{}
%\date{}

\begin{document}
	
	\maketitle
	

%----------------------------------------------------------------------
%----------------------------------------------------------------------
\begin{frame}[c]{Why Do We Need Evaluation in Machine Learning?}

\begin{itemize}
    \item Evaluation is the backbone of both machine learning and AutoML
    \item AutoML introduces additional stochasticity and selection effects that make na\"ive evaluation misleading
    \item Goal of this chapter: Get to know fundamental evaluation protocols and rigorous AutoML benchmarking practices
\end{itemize}
\end{frame}

\begin{frame}{Why is Evaluation Hard?}
    \begin{itemize}
        \item Data is finite, but we want to reason about infinite future data
        \item Models are stochastic and data is noisy
        \item Small design choices can change conclusions dramatically
        \item Strong optimization methods can exploit weaknesses in evaluation procedures resulting in (meta-)overfitting
    \end{itemize}
    \vspace{.5cm}
    \textbf{If evaluation is weak, progress becomes an illusion.}
\end{frame}


\begin{frame}{From Evaluation to Benchmarking}
\begin{itemize}
    \item Evaluation asks: How good is a model?
    \item Benchmarking asks: How good is a method compared to alternatives across many problems?
\end{itemize}
\vspace{0.5cm}
\begin{itemize}
    \item Benchmarks shape research directions
    \item Poor benchmarks reward shortcuts instead of real progress
    \item Good benchmarks enable reproducibility, comparability, scientific progress, and reusability
\end{itemize}
\end{frame}
    
\begin{frame}[c]{Why Is Evaluation Even More Important for AutoML?}
    \begin{itemize}
        \item AutoML systems typically explore many pipelines, hyperparameter settings, and thus models automatically
        \item The more configurations we try, the easier it becomes to overfit the evaluation process itself
        \item Automation amplifies both good and bad evaluation practices
    \end{itemize}
    \vspace{0.5cm}
    \textbf{AutoML applies evaluation methods systematically but validating them correctly is even more important}
\end{frame}

\begin{frame}{Guiding Questions for This Chapter}
\begin{itemize}
    \item How can we estimate future performance using limited data?
    \item When is one model truly better than another?
    \item When is an approach better than another?
    \item How do we avoid fooling ourselves when searching over thousands of models?
    \item How can we visualize optimization behavior rather than only final results?
\end{itemize}

\end{frame}

% \begin{frame}{Summary} 

% Most important insights from this introduction:
% \begin{enumerate}
%     \item Evaluation is our proxy for future performance
%     \item AutoML increases the need for careful evaluation because automation amplifies bias
%     \item Benchmarking turns evaluation into scientific evidence
% \end{enumerate}

% \vspace{0.5cm}

% $\leadsto$ The next parts of this chapter develop the tools needed to evaluate ML models and benchmark AutoML systems rigorously

% \end{frame}

\end{document}
